% !TEX program = xelatex
% =============================================================================
% cn-litigation · 劳动人事争议仲裁申请书
%
% 对标北京市劳动人事争议仲裁委员会官方表格文本（人社系统各地格式大同小异：
% 当事人信息表＋请求事项＋事实和理由＋签名日期＋附证据材料清单＋填表注意
% 事项）。官方原表为申请人/被申请人左右双栏，此处按本模板族统一的
% 左项目右内容要素表改排，要素项照录。
%
% 样张：虚构劳动争议，赵××（劳动者）申请示例商贸有限公司（用人单位）
% 支付拖欠工资与违法解除赔偿金；金额自洽——月工资 8,200 元，拖欠 3 个月
% =24,600 元；工作年限 2023-03-01 至 2026-05-31 为 3 年 3 个月，经济补偿
% 基数 3.5 个月，违法解除赔偿金按二倍即 7 个月×8,200=57,400 元。
%
% 编译：bash scripts/compile.sh templates/cn-litigation/labor-arbitration.tex --preview
% =============================================================================
\documentclass[zihao=-4]{ctexart}
\usepackage{litigation}

\begin{document}

\liTitle{劳动人事争议仲裁申请书}{}

\begin{liform}
\lisec{当事人信息}
\lirow{申请人\newline（劳动者）}{%
  \likv{姓名}{赵××}\likv{性别}{女}\likv{民族}{汉族}\likv{出生日期}{1990年9月15日}\newline
  \likv{身份证件类型}{居民身份证}\likv{证件号码}{110×××××××××××××××}\newline
  \likv{联系电话}{136\,0000\,0000}\likv{户口性质}{非农业}\newline
  \likv{住所地}{示例市示例区示例路88号}\newline
  \likv{通讯地址}{同住所地}}
\lirow{被申请人\newline（用人单位）}{%
  \likv{单位名称}{示例商贸有限公司}\likv{企业性质}{有限责任公司}\newline
  \likv{统一社会信用代码}{91××××××××××××××××}\newline
  \likv{法定代表人（负责人）}{钱××}\likv{职务}{执行董事}\newline
  \likv{单位联系人}{孙××}\likv{部门及职务}{人力资源部经理}\likv{联系电话}{010-8800\,0002}\newline
  \likv{住所地}{示例市示例区示例大街100号}\likv{办公地或经营地}{同住所地}\newline
  \likv{劳动合同履行地}{示例市示例区}\likv{是否签合同}{是}}
\lisec{请求事项}
\lirow{1.拖欠工资}{裁决被申请人支付申请人2026年3月1日至2026年5月31日期间拖欠的工资24,600元（月工资8,200元×3个月）}
\lirow{2.赔偿金}{裁决被申请人支付违法解除劳动合同赔偿金57,400元（申请人2023年3月1日入职，2026年5月31日被解除，工作年限3年3个月，经济补偿基数为3.5个月工资，赔偿金按二倍计7个月×8,200元）}
\lispan{其他补充：无}
\lisec{事实和理由}
\lispan{申请人于2023年3月1日入职被申请人处，任销售主管，双方订立书面劳动合同，约定月工资8,200元。2026年3月起，被申请人以经营困难为由停发工资。2026年5月31日，被申请人以“经营调整”为由口头通知解除劳动合同，未提前三十日书面通知，亦不符合《中华人民共和国劳动合同法》第三十九条、第四十条、第四十一条规定的任一解除情形，属违法解除。依据《劳动合同法》第三十条、第四十七条、第四十八条、第八十七条之规定，提出上述仲裁请求。}
\lisec{证据材料清单（另页附后）}
\lirow{1.劳动合同}{原件\ckboxv 复印件\ckbox\quad 证明：双方存在劳动关系及工资标准}
\lirow{2.银行工资流水}{原件\ckboxv 复印件\ckbox\quad 证明：月工资数额及2026年3月起停发事实}
\lirow{3.解除通知录音及考勤记录}{原件\ckboxv 复印件\ckbox\quad 证明：被申请人违法解除劳动合同}
\end{liform}

\liSubmitTo{示例市劳动人事争议仲裁委员会}

\liSignLine{申请人（签名或盖章）：}{}
\liSign{2026年\hspace{1\ccwd}6月\hspace{1\ccwd}20日}

\begin{linote}
注：1.本申请书须用钢笔或签字笔填写，一式三份，其中两份交劳动人事争议仲裁委员会，一份申请人自存；被申请人为两个或以上的，申请书须相应增加。\newline
\hspace*{2\ccwd}2.申请内容填写不下时，可在申请书中间加页。\newline
\hspace*{2\ccwd}3.填写的申请书须字迹清楚，文字简练规范。\newline
\hspace*{2\ccwd}4.证据材料清单及相应证据材料另页附后，一式三份，其中一份为原件由申请人自存，另两份为复印件交劳动人事争议仲裁委员会。\newline
\hspace*{2\ccwd}5.本申请书（含证据材料清单及相应证据材料）及其复印件均需使用A4型纸。
\end{linote}

\end{document}
